\documentclass[12pt, twoside]{article}
\input{/Users/chris/github/course-files/docs/ib/preamble}

\fancyhead[LE]{\thepage}
\fancyhead[RO]{Name: \hspace{4cm} \,\\ \hspace{2.9cm} \,}
\fancyhead[LO]{La Scuola d'Italia / Dr.\ Huson / IB Math HL \\* 21 May 2026}

\begin{document}
\subsubsection*{5.1 Quiz: Maclaurin Series (no calculator)}

Answer all three problems. Show full working in the spaces provided.
Total: 28 marks. Time: 45 minutes.

\begin{enumerate}[itemsep=0.15cm]

%--- Problem 1: Maclaurin of (1-x)^-4 with depreciation application [7 marks] ---
%--- Source: AA HL 2025 Nov P1 TZ1 Q05 ---
\item[1.] \makebox[\linewidth][s]{[Maximum mark: 7]\hfill {\small AA HL 2025 Nov P1 TZ1 Q05}}

The first four terms of the Maclaurin series expansion of $(1 - x)^{-4}$ are
\[
1 + ax + bx^2 + 20x^3,
\]
where $a, b \in \mathbb{Z}^{+}$.

\begin{enumerate}[itemsep=0.15cm]
  \item \begin{enumerate}[itemsep=0.15cm]
    \item[(i)] Show that $a = 4$. \hfill [2]
    \item[(ii)] Find the value of $b$. \hfill [3]
  \end{enumerate}
  \item A car was purchased four years ago. The car depreciated in value by
  $10\%$ each year. The value of the car today is $\$1000$. Using the results
  of part (a), estimate the value of the car four years ago. \hfill [2]
\end{enumerate}

\begin{tikzpicture}
  \draw (-0.6,0) rectangle (11,11);
  \foreach \y in {1,2,...,10} \draw[dotted] (0,\y)--(10.5,\y);
\end{tikzpicture}

\newpage

%--- Problem 2: Maclaurin of sin(x^2) and product manipulation [7 marks] ---
%--- Source: AA HL 2024 May P1 TZ1 Q08 ---
\item[2.] \makebox[\linewidth][s]{[Maximum mark: 7]\hfill {\small AA HL 2024 May P1 TZ1 Q08}}

\begin{enumerate}[itemsep=0.15cm]
  \item Find the first two non-zero terms in the Maclaurin series of
  \begin{enumerate}[itemsep=0.15cm]
    \item[(i)] $\sin\!\left(x^{2}\right)$; \hfill [2]
    \item[(ii)] $\sin^{2}\!\left(x^{2}\right)$. \hfill [3]
  \end{enumerate}
  \item Hence, or otherwise, find the first two non-zero terms in the
  Maclaurin series of $4x \sin\!\left(x^{2}\right) \cos\!\left(x^{2}\right)$.
  \hfill [2]
\end{enumerate}

\begin{tikzpicture}
  \draw (-0.6,0) rectangle (11,11);
  \foreach \y in {1,2,...,10} \draw[dotted] (0,\y)--(10.5,\y);
\end{tikzpicture}

\newpage

%--- Problem 3: Induction + Maclaurin + limit [14 marks] ---
%--- Source: AA HL 2021 Nov P1 TZ0 Q11 ---
\item[3.] \makebox[\linewidth][s]{[Maximum mark: 14]\hfill {\small AA HL 2021 Nov P1 TZ0 Q11}}

\begin{enumerate}[itemsep=0.15cm]
  \item Prove by mathematical induction that
  \[
  \frac{d^{n}}{dx^{n}}\!\left(x^{2} e^{x}\right)
  = \left[x^{2} + 2nx + n(n - 1)\right] e^{x}
  \quad\text{for } n \in \mathbb{Z}^{+}.
  \]
  \hfill [7]
\end{enumerate}

\begin{tikzpicture}
  \draw (-0.6,0) rectangle (11,17);
  \foreach \y in {1,2,...,16} \draw[dotted] (0,\y)--(10.5,\y);
\end{tikzpicture}

\newpage

\begin{enumerate}[itemsep=0.15cm]
  \item[(b)] Hence or otherwise, determine the Maclaurin series of
  $f(x) = x^{2} e^{x}$ in ascending powers of $x$, up to and including
  the term in $x^{4}$. \hfill [3]
\end{enumerate}

\begin{tikzpicture}
  \draw (-0.6,0) rectangle (11,9);
  \foreach \y in {1,2,...,8} \draw[dotted] (0,\y)--(10.5,\y);
\end{tikzpicture}

\begin{enumerate}[itemsep=0.15cm]
  \item[(c)] Hence or otherwise, determine the value of
  \[
  \lim_{x \to 0}\!\left[\frac{\left(x^{2} e^{x} - x^{2}\right)^{3}}{x^{9}}\right].
  \]
  \hfill [4]
\end{enumerate}

\begin{tikzpicture}
  \draw (-0.6,0) rectangle (11,9);
  \foreach \y in {1,2,...,8} \draw[dotted] (0,\y)--(10.5,\y);
\end{tikzpicture}

\end{enumerate}
\end{document}
